% Results-neutral draft for the final paper.
% Do not add outcome claims until the 32-cell grid, frozen-threshold scoring,
% and once-only 50-scene verified evaluation are complete. The submission
% version should report measured low-label ordering, label savings, and
% architecture generality without seed-variance claims.
\documentclass[10pt]{article}

\usepackage[T1]{fontenc}
\usepackage[letterpaper,margin=1in]{geometry}
\usepackage{microtype}
\usepackage[hidelinks]{hyperref}

\pagestyle{empty}

\title{Label-Efficient Dark-Vessel Detection in SAR:\\
Does SAR-Domain Pretraining Beat Optical and ImageNet Transfer Across ViT and CNN?}
\author{John Roth \and Kyle Wagner\\
Johns Hopkins University}
\date{}

\begin{document}
\maketitle

\begin{abstract}
Detecting dark vessels---ships without an Automatic Identification System
correlate---in spaceborne synthetic aperture radar (SAR) is operationally
important, but human-verified examples are scarce. We investigate whether
pretraining domain improves label efficiency on xView3-SAR and whether any
benefit generalizes across architecture families. Our controlled design uses
two size-matched tracks: ViT-B/16 (approximately 86M parameters) and
ConvNeXt-V2-Base (approximately 89M). Within each track, we compare random
initialization with downloaded encoders trained on generic ImageNet, optical
remote-sensing imagery, or SAR imagery. All eight arms share a fixed
three-channel $[\mathrm{VH},\mathrm{VV},\mathrm{VH}{-}\mathrm{VV}]$ input,
point-native heatmap detector, optimizer, schedule, scene-level splits, and
seed, and are fine-tuned at nested 10\%, 25\%, 50\%, and 100\% training-scene
fractions. We measure distance-matched vessel-detection F1 on held-out scenes
and, on a disjoint once-touched set of 50 human-verified scenes, dark-vessel
recall and near-shore F1. All values are reported as seed-0 point estimates;
external YOLO and zero-shot vision-language references remain separate from
the controlled curves. By separating initialization comparisons within an
architecture from matched-role comparisons across architectures, the study
provides a reproducible test of when pretrained representations reduce
labeling demand for maritime SAR detection.
\end{abstract}

\end{document}
